\section{Setting}
\label{sec:setting}

\subsection{Models as maps}
\label{sec:models}

All models considered here are autonomous maps that are iterated to produce a trajectory. An echo state
network\cite{Jaeger2001EchoState,LukoseviciusJaeger2009Reservoir} has the reservoir update
$r_{k+1}=\tanh(Wr_k+W_{\mathrm{in}}u_k+b)$ with a fixed random matrix $W$, input matrix $W_{\mathrm{in}}$ and bias $b$, and
a linear readout $u_k=W_{\mathrm{out}}r_k$ fitted by ridge regression. In closed loop the output is fed back as the input,
and the model is the map
\begin{equation}
F(r)=\tanh(Ar+b),\qquad A=W+W_{\mathrm{in}}W_{\mathrm{out}},
\label{eq:esn}
\end{equation}
on $\mathbb{R}^N$, with $N$ between $50$ and $2000$ for the trained networks. The other three classes act directly on the
normalized physical state
$u=(x-\bar x)/\sigma$, with the mean $\bar x$ and the standard deviation $\sigma$ of the training data taken per
coordinate: the one-hidden-layer perceptron (``mlp'') $F(u)=W_2\tanh(W_1u+b_1)+b_2$, the two-hidden-layer
perceptron (``mlp2'') $F(u)=W_3\tanh\bigl(W_2\tanh(W_1u+b_1)+b_2\bigr)+b_3$, and the residual network (``res'')
$F(u)=u+W_2\tanh(W_1u+b_1)+b_2$, a forward-Euler step of a neural vector
field.\cite{He2016ResNet,E2017Proposal,HaberRuthotto2018Stable} The first three classes have bounded output; the residual
network does not, and its trajectories can escape to infinity as those of the true systems do.

The equilibria of a model are the fixed points of its map. An equilibrium is unstable if $DF$ has an eigenvalue of
modulus greater than $1$. We write $U$ for the number of unstable equilibria found and $\Uin$ for the number of those whose
physical image ($W_{\mathrm{out}}r^*$ for the echo state network, $\sigma u^*+\bar x$ otherwise) lies in the region $Q$ of the
true system (Table~\ref{tab:systems}). For the perceptrons almost all unstable equilibria lie in $Q$, while residual networks
have on average one or two more outside it (Sec.~\ref{sec:classes}); for echo state networks we count all equilibria and
write $U$.

\subsection{Hidden attractors of a model}
\label{sec:hidden}

An attractor $\mathcal{A}$ of a map or flow, with basin $B$, is \emph{excited} from an equilibrium $e$ if $B$ intersects
every neighborhood of $e$. It is \emph{self-excited} if it is excited from at least one equilibrium and \emph{hidden}
otherwise.\cite{LeonovKuznetsov2013,Dudkowski2016} The definition is the same for maps, which have hidden attractors of
their own;\cite{Jiang2016Maps} a trained model is a map even when the true system is a flow. We consider attractors with an open basin (local attractors
in the sense of Leonov and Kuznetsov), which includes the chaotic attractors and limit cycles of the true systems; for them
an intersection of $B$ with a neighborhood of $e$ is open and has positive measure, which is what a numerical test can
detect. Attractors in Milnor's sense,\cite{Milnor1985Attractor} whose basin need not be open, would need a measure-theoretic
version of the definition. Hiddenness is a property of the pair
(attractor, set of equilibria), so the same attractor can be hidden in the true system and self-excited in a model that
has different equilibria.

A model \emph{reproduces} a hidden attractor of the true system if trajectories started from states on the true attractor
settle on a model attractor of the same type (chaotic or regular) that is close to it in two senses: the normalized
distance of first and absolute moments, $\Delta_{\mathrm{att}}$, and the Hausdorff distance relative to the attractor
diameter, $d_H$, are both below $0.2$ (Appendix~\ref{app:protocol}). The question of this paper is whether a reproduced
attractor stays hidden, now with respect to the equilibria of the model. We call the fraction of reproductions in which it
does the \emph{survival} probability.

\subsection{True systems}
\label{sec:systems}

\begin{table*}
\caption{Test systems. Equilibria are listed with their flow-map index $\iota$ in parentheses, Eq.~\eqref{eq:trueindex};
$s$ is the index sum. $p_B$ is the fraction of uniform initial conditions in the region $Q$ that reach the hidden
attractors (both of them for Chua's circuit); $\Delta t$ is the sampling interval of the training series. Equations and
regions are given in Appendix~\ref{app:systems}. $^{a}$Hidden only in a weak sense: not reached from the stable half-line of
equilibria, but reached from the segment $0.35\lesssim z\lesssim2.2$ of the unstable one, hence self-excited in the strict
sense; reported separately (Sec.~\ref{sec:systems}).}
\label{tab:systems}
\begin{ruledtabular}
\begin{tabular}{lllllllll}
Regime & System & Parameters & Equilibria ($\iota$) & $s$ & Hidden & Visible & $p_B$ & $\Delta t$ \\
\hline
SE3 & one stable equilibrium & -- & stable ($+$) & $1$ & chaotic & equilibrium & 0.152 & 0.2 \\
SE11 & one stable equilibrium & -- & stable ($+$) & $1$ & chaotic & equilibrium & 0.00827 & 0.15 \\
LE1 & line of equilibria & $a=15$, $b=1$ & line $x=y=0$ (stable for $z<0$) & -- & chaotic$^{a}$ & equilibria & 0.00127 & 0.16 \\
WC & Wang--Chen & $a=0.006$ & stable ($+$) & $1$ & chaotic & equilibrium & 0.67 & 0.18 \\
WC$'$ & Wang--Chen & $a=0.03$ & stable ($+$) & $1$ & cycle & equilibrium & 0.668 & 0.18 \\
Chua I & Chua's circuit & set I & 2$\times$ saddle-focus ($-$), stable ($+$) & $-1$ & 2$\times$ chaotic & equilibrium & 0.0369 & 0.07 \\
Chua II & Chua's circuit & set II & saddle-focus ($-$), 2$\times$ stable ($+$) & $1$ & 2$\times$ chaotic & cycle & 0.00274 & 0.05 \\
\end{tabular}
\end{ruledtabular}
\end{table*}


Seven regimes of five three-dimensional flows with attractors reported as hidden and at least one equilibrium make up the
test set (Table~\ref{tab:systems}); six of them are hidden in the strict sense.\cite{Molaie2013,JafariSprott2013,WangChen2012,Leonov2011PLA,Stankevich2017} Two systems from
the family with one stable equilibrium (SE3 and SE11) have a hidden chaotic attractor next to a stable focus. LE1 has a line
of equilibria $x=y=0$, stable for $z<0$ and unstable for $z>0$. Its attractor is hidden only in the weaker sense of
Jafari and Sprott,\cite{JafariSprott2013} who noted that its basin meets the line in some portions: it is not reached from
the stable half-line, but it is reached from a segment of the unstable one, $0.35\lesssim z\lesssim2.2$: trajectories
started at distance $10^{-8}$ from the line reach it along $43\%$ of the directions on average for $0.3<z\le1$, and along
all of them at some points (Appendix~\ref{app:systems}).
In the sense of Sec.~\ref{sec:hidden} the attractor of LE1 is therefore self-excited. We report LE1 separately, as a
control in which the true system itself has exciting equilibria: for it, an attractor that is not excited in a model
marks a loss of faithfulness rather than the survival of hiddenness. All pooled statistics use the other six regimes.
The Wang--Chen system has a single stable equilibrium and a hidden chaotic attractor at
$a=0.006$ that becomes a hidden limit cycle at $a=0.03$; its basin covers about two thirds of the region. Chua's circuit is
taken in the two classical regimes: in regime~I a stable origin and two saddle-foci coexist with two symmetric hidden
chaotic attractors; in regime~II two stable equilibria and an unstable origin coexist with a large visible limit cycle
and two hidden chaotic attractors. The equilibria of LE1 are not isolated and have no index, so
Proposition~\ref{prop:forced} applies to the other six regimes; of these, only Chua~I has $s\neq1$. As a control with $s=0$ we
use the four-dimensional Li--Sprott system, which has no equilibria, in the setting of R\"ohm
\textit{et al.}:\cite{Rohm2021} networks trained on one of its attractors reproduce coexisting attractors that were not in
the data (Appendix~\ref{app:systems}).

\subsection{Training data}
\label{sec:data}

Every model is trained on a noisy sampled time series of the true system. Three training sets differ in what they show
the model. \emph{H1} is one trajectory of $10^4$ samples on the hidden attractor. \emph{E1} consists of $20$ short
trajectories, $500$ samples each, started in the basin of the visible attractor or equilibrium; they show the transients
toward it and nothing of the hidden attractor except where these transients pass close to it. \emph{HE1} is the union of
H1 and E1. For the intervention experiment of Sec.~\ref{sec:data_effects} the reservoir of an echo state network is kept
fixed and only the readout is refitted on H1 plus $5$, $10$ or $20$ transients (E5, E10, E20). Integration steps, sampling
intervals, noise levels and regions are listed in Appendix~\ref{app:systems}.
